\documentclass[10pt,a4paper]{amsart}
\usepackage[margin=19mm]{geometry}
\usepackage{amsmath,amssymb,mathtools}
\usepackage[scr=rsfs,cal=euler]{mathalfa}
\usepackage{xfrac}
\usepackage[unicode,hidelinks,bookmarks=false]{hyperref}
\newcommand{\ie}{i.e.\ }
\newcommand{\cf}{cf.\ }
\newcommand{\resp}{respectively\ }
\newcommand{\Subsection}{\subsection}
\providecommand{\nobreakdash}{\nobreak\hbox{-}}
\setlength{\emergencystretch}{2em}
\multlinegap0pt
\usepackage{bm}
\usepackage{bbm}
\usepackage{dsfont}
\usepackage{xspace}
\usepackage{cancel}
\usepackage{mathtools}
\usepackage{amsthm}
\makeatletter
\@ifundefined{alphthm}{\newtheorem{alphthm}{Theorem}}{}
\@ifundefined{theorem}{\newtheorem{theorem}{Theorem}}{}
\@ifundefined{claim}{\newtheorem{claim}[theorem]{Claim}}{}
\@ifundefined{corollary}{\newtheorem{corollary}[theorem]{Corollary}}{}
\@ifundefined{definition}{\newtheorem{definition}[theorem]{Definition}}{}
\@ifundefined{lemma}{\newtheorem{lemma}[theorem]{Lemma}}{}
\@ifundefined{notation}{\newtheorem{notation}[theorem]{Notation}}{}
\@ifundefined{proposition}{\newtheorem{proposition}[theorem]{Proposition}}{}
\makeatother
\newcommand{\acl}{\operatorname{acl}}
\title[Topological reconstruction theorems over uncountable algebra]{Topological reconstruction theorems over uncountable algebraically closed fields}
\author{Benjamin Castle, Ronan O'Gorman}
\date{Source: arXiv:2607.14472v1 (CC BY 4.0)}
\begin{document}
\maketitle

\noindent\emph{Source and licence.} Topological reconstruction theorems over uncountable algebraically closed fields. Version arXiv:2607.14472v1 (CC BY 4.0), distributed under the Creative Commons Attribution 4.0 International licence, as recorded in the arXiv licence metadata for this exact version and shown on the abstract page https://arxiv.org/abs/2607.14472v1. Full rights notice: licenses/arxiv-2607.14472-CC-BY-4.0.txt.\par\smallskip

\noindent\textit{Editorial scope.} Counted unit: the Proposition whose source label is \texttt{P: preservation inductive step} in the article (source0.tex lines 1076--1078, source label \texttt{P: preservation inductive step}), delivered together with the article's own complete original proof of it (source0.tex lines 1079--1145, one \texttt{proof} environment of 67 lines). No mathematics is written, repaired or re-derived: statement and proof are the article's own text line for line. Required context retained uncounted: T: intro link preservation (lines 361--362); L: codim 1 definable over acl (lines 472--474); L: good preservation (lines 675--682); D: generic sweeping witness (lines 701--709); L: generic sweeping witness (lines 711--713); L: relativized sweeping is unique (lines 739--741); C: non-generics are small (lines 942--947); L: linked restriction (lines 1006--1008); T: preservation (lines 1017--1023). Every definition, display and label of those lines is retained unchanged. Omitted from this package and not redistributed: 1--360, 363--471, 475--674, 683--700, 710--710, 714--738, 742--941, 948--1005, 1009--1016, 1024--1075, 1146--2215. Nothing inside a retained excerpt is omitted or altered: every retained body is delivered exactly as the article's lines are written, so no omission receipt of any kind accompanies this package. Citations: the retained excerpts carry no citation command at all, so no bibliography entry of the article is transcribed and no reference is redistributed. Reference adaptation: none was needed and none was made; every reference inside the delivered unit resolves inside it, so no unresolved reference or citation is produced. Printed slips kept literal: no printed word, symbol, hypothesis or number of the counted statement or of its proof was altered. Layout and class: the article's own class is \texttt{amsart} with packages including amssymb, hyperref, graphicx, float, xy, comment, tikz; the wrapper is the lane's pinned \texttt{amsart} editorial wrapper at 10pt on A4, which restates the theorem-like environments the retained text needs and adds the article's own macro definitions that the retained text needs (\texttt{\textbackslash acl}), with extra wrapper packages (bm, bbm, dsfont, xspace, cancel, mathtools). The delivered numbering is the unit's own. Assets: no retained span contains a figure environment, a graphics-inclusion command, a PDF-page inclusion, a table or a TikZ picture; no image file is copied into this package, the package declares no asset dependency and no image file is redistributed with it. Licence: version arXiv:2607.14472v1 of this article is distributed under the Creative Commons Attribution 4.0 International licence, as recorded in the arXiv licence metadata for this exact version and shown on the abstract page https://arxiv.org/abs/2607.14472v1. A full rights notice is delivered at licenses/arxiv-2607.14472-CC-BY-4.0.txt. Characters: every retained excerpt is pure ASCII, so the delivered engine, XeLaTeX with its default Latin Modern text font, reports no missing character.
\begin{alphthm}[= Theorem \ref{T: preservation}]\label{T: intro link preservation} Fix $O_2$ as above. Then for all but finitely many closed irreducible codimension 1 subvarieties $Y\subset X_1$, we have that $O_1$ $n$-sweeps $Y_1$ if and only if $O_2$ $n$-sweeps $\varphi(Y_1)$.
\end{alphthm}

\begin{lemma}\label{L: codim 1 definable over acl}
    Let $W$ be an $A$-definable irreducible quasi-projective variety over an uncountable algebraically closed field, and let $V\subset W$ be a closed irreducible codimension 1 subvariety defined over $B\supset A$. Suppose that some $B$-generic element of $V$ is not $A$-generic in $W$. Then $V$ is defined over $\acl(A)$.
\end{lemma}

\begin{lemma}\label{L: good preservation}
    Let $O\in\operatorname{Orb}(A)$ be uniquely $k$-sweeping.
    \begin{enumerate}
        \item $X$ is $O$-good.
        \item Suppose $k\geq 2$ and $O$ uniquely sweeps $(Z_1,...,Z_k)$, where each $Z_i$ is $A$-definable. Let $Y\subset X$ be $A$-definable and $O$-good. Then for $A$-generic $a\in Z_1$, $Y$ remains $O\restriction a$-good. Moreover, for $A$-generic $b\in Z_2\times...\times Z_k$, $Y$ also remains $O\restriction b$-good.
        
    \end{enumerate}
\end{lemma}

\begin{definition}\label{D: generic sweeping witness}
    Suppose $O\in\operatorname{Orb}(A)$ is irreducible, $Z_1,...,Z_k$ are $A$-definable closed irreducible subvarieties, and $O$ sweeps $(Z_1,...,Z_k)$. A \textit{generic witness to $O$ sweeping $(Z_1,...,Z_k)$} consists of a member $Y\in O$, with canonical base $c$, and a tuple $a=(a_1,...,a_k)\in Z_1\times...\times Z_k$, so that:
    \begin{enumerate}
        \item $a\in Y^k$.
    \item $a$ is $A$-generic in $Z_1\times...\times Z_k$.
    \item $\dim(Y\cap Z_i)=\dim(Z_i)-1$ for all $i$.
    \item $a$ is $Ac$-generic in $(Y\cap Z_1)\times...\times(Y\cap Z_k)$.
    \end{enumerate}
\end{definition}

\begin{lemma}\label{L: generic sweeping witness}
   Let $O,A,Z_1,...,Z_k$ be as in Definition \ref{D: generic sweeping witness}. If each $Z_i$ is $O$-good, then there is a generic witness to $O$ sweeping $(Z_1,...,Z_k)$.
\end{lemma}

\begin{lemma}\label{L: relativized sweeping is unique}
    Let $O\in\operatorname{Orb}(A)$ be uniquely $k$-sweeping, and let $Z_1,...,Z_k$ be $A$-definable $O$-good subvarieties. If $O$ sweeps $(Z_1,...,Z_k)$, then $O$ uniquely sweeps $(Z_1,...,Z_k)$.
\end{lemma}

\begin{corollary}\label{C: non-generics are small}
    \begin{enumerate}
        \item For any $(A_1,A_2)$, the set $N_k$ of $a\in X_1^k$ which are not $(A_1,A_2)$-generic is small.
        \item If $S\subset X_1^k$ is large, then for any $(A_1,A_2)$, the set $S'$ of $(A_1,A_2)$-generics in $X_1^k$ which belong to $S$ is also large.
    \end{enumerate}
\end{corollary}

\begin{lemma}\label{L: linked restriction}
    Suppose $k>1$ and $O_i\in\operatorname{Orb}(A_i)$ are uniquely $k$-sweeping and linked. Then there is a large set $S\subset X_1$ consisting of $(A_1,A_2)$-generic elements $a\in X_1$ so that $O_1\restriction a$ and $O_2\restriction\varphi(a)$ are linked. 
\end{lemma}

\begin{theorem}[= Theorem \ref{T: intro link preservation}]\label{T: preservation}
    Let $O_1,O_2$ be linked, where each $O_i$ is a uniquely $k$-sweeping $A_i$-orbit in $X_i$. Suppose $\varphi(Z_1^j)=Z_2^j$ for $j=1,...,k$, where each $Z_i^j$ is $O_i$-good. Then:
    \begin{enumerate}
        \item $O_1$ sweeps $(Z_1^1,...,Z_1^k)$ if and only if $O_2$ sweeps $(Z_2^1,...,Z_2^k)$.
        \item Suppose each $O_i$ $k$-sweeps $(Z_i^1,...,Z_i^k)$, and each $Z_i^j$ is $A_i$-definable. Then for any $(A_1,A_2)$-generic $(a_1,a_2)\in(Z_1^1\times...\times Z_1^k,Z_2^1\times...\times Z_2^k)$ we have $\varphi(O_1(a_1))=O_2(a_2)$.
    \end{enumerate}
\end{theorem}

\begin{proposition}\label{P: preservation inductive step}
    Let $k,O_i,A_i,Z_i^j$ be as in Theorem \ref{T: preservation}. Assume that $k\geq 2$ and Theorem \ref{T: preservation} holds for $k'<k$. Then Theorem \ref{T: preservation} holds for $k,O_i,A_i,Z_i^j$.
\end{proposition}

\begin{proof}
    As in the base case, we assume throughout that each $Z_i^j$ is $A_i$-definable. Throughout the proof, fix a large set $S\subset X_1$ so that for all $a\in S$ the orbits $O_1\restriction a$ and $O_2\restriction\varphi(a)$ are linked (which exists by Lemma \ref{L: linked restriction}). Throughout the proof, we will apply the inductive hypothesis to different restrictions of the $O_i$ at generic points. Each time, we are using that the $Z_i^j$ remain good after generic restrictions (i.e. Lemma \ref{L: good preservation}(2)).
    
    Assume that $O_1$ sweeps $(Z_1^1,...,Z_1^k)$. We will show that $O_2$ sweeps $(Z_2^1,...,Z_2^k)$, and that for all $(A_1,A_2)$-generic elements $(a_1,a_2)\in(Z_1^1\times...\times Z_1^k,Z_2^1\times...\times Z_2^k)$ we have $\varphi(O_1(a_1))=O_2(a_2)$. By symmetry, this is enough to prove the whole proposition.

    We proceed with a sequence of claims.

    \begin{claim}\label{Cl: preservation inductive step claim 1} For generic $a\in X_1$, $O_1\restriction a_1$ sweeps $(Z_1^2,...,Z_1^k)$.
    \end{claim}
    \begin{proof}
        Let $z_1=(z_1^1,...,z_1^k)\in Z_1^1\times...\times Z_1^k$ be $A_1$-generic, and let $Y_1=O_1(z_1)$. If some generic element $y_1\in Y_1$ is $A_1z_1^2...z_1^k$-generic in $X_1$, then $Y_1$ belongs to $O(y_1)$ and contains the $A_1z_1^2...z_1^k$-generic point $(z_1^2,...,z_1^k)\in Z_1^2\times...\times Z_1^k$, which implies the claim. Otherwise, assume every generic element of $Y_1$ is not $A_1z_1^2...z_1^k$-generic in $X_1$. Then by Lemma \ref{L: codim 1 definable over acl}, $Y_1$ is $\operatorname{acl}(A_1z_1^2...z_1^k)$-definable. But then $$\dim(z_1^1/A_1z_1^2...z_1^k)\leq\dim(Y_1\cap Z_1^1)<\dim(Z_1^1),$$ contradicting the genericity of $z$ in $Z_1^1\times...\times Z_1^k$. (Note that the inequality $\dim(Y_1\cap Z_1^1)<\dim(Z_1^1)$ is because $Z_1^1$ is $O_1$-good, thus not $O_1$-common). 
    \end{proof}

    \begin{claim}\label{Cl: claim 2 inductive step}
        For generic $a_2\in X_2$, $O_2\restriction a_2$ sweeps $(Z_2^2,...,Z_2^k)$.
    \end{claim}
    \begin{proof}
        Let $a_1\in S$ be $(A_1,A_2)$-generic. By Claim \ref{Cl: preservation inductive step claim 1}, $O_1\restriction a_1$ sweeps $(Z_1^2,...,Z_1^k)$. So by the inductive hypothesis, $O_2\restriction a_2$ sweeps $(Z_2^2,...,Z_2^k)$. By automorphism invariance, the same holds for all $A_2$-generic elements of $X_2$.
    \end{proof}

    \begin{claim}\label{Cl: claim 3 inductive step}
    Each $O_i$ sweeps $(X_i,Z_i^2,...,Z_i^k)$.
    \end{claim}
    \begin{proof}
        Clear be the previous two claims. That is, let $a\in X_i$ be $A_i$-generic. We showed above that $O_i\restriction a$ sweeps $(Z_i^2,...,Z_i^k)$, so there are $Y\in O_i\restriction a$ and an $A_ia$-generic $b\in Z_i^2\times...\times Z_i^k$ with $b\in Y^{k-1}$. The $(a,b)$ is $A_i$-generic in $X_i\times Z_i^2\times...\times Z_i^k$ and contained in $Y\in O_i$, which gives the claim.
    \end{proof}

    By Claim \ref{Cl: claim 3 inductive step}, Lemma \ref{L: good preservation}, and Lemma \ref{L: relativized sweeping is unique}, each $O_i$ uniquely sweeps $(X_i,Z_i^2,...,Z_i^k)$. Thus, if $a\in X_i$ and $b\in Z_i^2\times...\times Z_i^k$ are independent generics over $A_i$, then $O_i(a,b)\in O_i$ is well-defined. Now recall:

    \begin{notation} For $A_i$-generic $b\in Z_i^2\times...\times Z_i^k$, $O_i\restriction b$ denotes the $A_ib$-orbit consisting of all members of $O_i$ of the form $O_i(a,b)$ for $A_ib$-generic $a\in X_i$.
    \end{notation}

    Since $O_i$ uniquely sweeps $(X_i,Z_i^2,...,Z_i^k)$, it is clear that for $A_i$-generic $b\in Z_i^2\times...\times Z_i^k$ the orbit $O_i\restriction b$ is uniquely 1-sweeping.

    \begin{claim}\label{Cl: claim 4 inductive step}
        Let $(b_1,b_2)\in(Z_1^2\times...\times Z_1^k,Z_2^2\times...\times Z_2^k)$ be $(A_1,A_2)$-generic. Then $O_1\restriction b_1$ and $O_2\restriction b_2$ are linked.
    \end{claim}
    \begin{proof}
        Recall that $S\subset X_1$ is a large set so that $O_1\restriction a$ and $O_2\restriction\varphi(a)$ are linked for all $a\in S$. By Corollary \ref{C: non-generics are small}(2), there is a large set $S'\subset S$ whose elements are all $(A_1b_1,A_2b_2)$-generic.

        We show that for all $(a_1,a_2)$ with $a_1\in S'$ and $\varphi(a_1)=a_2$ we have $(O_1\restriction b_1)(a_1)=(O_2\restriction b_2)(a_2)$. Since $S'$ is large, this gives the claim. So, let $\varphi(a_1)=a_2$ with $a_1\in S'$. Note for each $i$ that $$(O_i\restriction b_i)(a_i)=O_i(a_i,b_i)=(O_i\restriction a_i)(b_i).$$ So we can equivalently show that $\varphi$ sends $(O_1\restriction a_1)(b_1)$ to $(O_2\restriction a_2)(b_2)$. But this is automatic by the inductive hypothesis of Theorem \ref{T: preservation} (applied to the orbits $O_i\restriction a_i$, which sweep $(Z_i^2,...,Z_i^k)$).
    \end{proof}

    \begin{claim}\label{Cl: claim 5 inductive step}
    Let $(b_1,b_2)\in(Z_1^2\times...\times Z_1^k,Z_2^2\times...\times Z_2^k)$ be $(A_1,A_2)$-generic. Then each $O_i\restriction b_i$ sweeps $Z_i^1$.
    \end{claim}
    \begin{proof}
        By Claim \ref{Cl: claim 4 inductive step} and the base case (applied to $O_i\restriction b_i$), it suffices to show this for $i=1$. But the case $i=1$ follows since $O_1$ sweeps $(Z_1^1,...,Z_1^k)$. Namely, let $z\in Z_1^1$ be $A_1b_1$-generic. Then by Lemma \ref{L: generic sweeping witness}, there is $Y\in O_1$ so that $(z,b_1)\in Y^k$ is a generic witness to $O_1$ sweeping $(Z_1^1,...,Z_1^k)$. Let $a\in Y$ be generic over $A_1zb_1$. If $a$ is not $A_1b_1$-generic in $X_1$, then by Lemma \ref{L: codim 1 definable over acl}, $Y$ is $\acl(A_1b_1)$-definable; but then since $z$ is $A_1b_1$-generic in $Z_1^1$, we get $Z_1^1\subset Y$, contradicting that $Z_1^1$ is not $O_1$-common. So $a$ must be $A_1b_1$-generic in $X_1$. But then $(a,b_1)\in X_1\times Z_1^2\times...\times Z_1^k$ is $A_1$-generic, so that $Y=O_1(a,b_1)$, and thus $Y\in O_1\restriction b_1$. So there is a member of $O_1\restriction b_1$ containing the $A_1b_1$-generic element $z\in Z_1^1$, which gives the claim. 
    \end{proof} 

    Finally, we are ready to prove the two clauses of Theorem \ref{T: preservation}.

    \begin{claim}
        $O_2$ sweeps $(Z_2^1,...,Z_2^k)$.
    \end{claim}
    \begin{proof}
       Let $(b_1,b_2)\in(Z_1^2\times...\times Z_1^k,Z_2^2\times...\times Z_2^k)$ be $(A_1,A_2)$-generic, and let $z\in Z_2^1$ be $A_2b_2$-generic. By Claim \ref{Cl: claim 5 inductive step}, there is $Y\in O_2\restriction b_2$ with $z\in Y$. Then $(z,b_2)\in Z_2^1\times...\times Z_2^k$ is $A_2$-generic and contained in $Y\in O_2$, as desired.
    \end{proof}

    \begin{claim} Let $(x_1,x_2)\in(Z_1^1\times...\times Z_1^k,Z_2^1\times...\times Z_2^k)$ be $(A_1,A_2)$-generic. Then $\varphi(O_1(x_1))=O_2(x_2)$.
    \end{claim}
    \begin{proof}
        Write $x_i=(z_i,b_i)$ with $z_i\in Z_i^1$ and $b_i\in Z_i^2\times...\times Z_i^k$. Then $$\varphi(O_1(x_1))=\varphi((O_1\restriction b_1)(z_1))=(O_2\restriction b_2)(z_2)=O_2(x_2),$$ where the middle equality is by the inductive hypothesis.
    \end{proof}

    Finally, by the two most recent claims, the proof of Theorem \ref{T: preservation} is complete.
\end{proof}

\end{document}
